\chapter{Neonatal Jaundice / Hyperbilirubinemia / Kernicterus}
\index{Neonatal Jaundice / Hyperbilirubinemia / Kernicterus}
\label{sec:Neonatal Jaundice}

\section{Overview}
Neonatal jaundice (hyperbilirubinemia) is the yellow discoloration of the skin and sclerae due to elevated serum bilirubin, occurring in 60--80\% of term and nearly all preterm newborns. This pregnancy-related condition requires careful monitoring to ensure the safety of both mother and baby. Early detection and management are essential.
\section{Causes}
The underlying mechanisms involve complex interactions between genetic predisposition and environmental factors. Disruption of normal physiological processes leads to the characteristic features of the condition. Multiple contributing factors may act synergistically to trigger or worsen the condition.

\begin{itemize}[noitemsep]
  \item Physiologic jaundice results from increased bilirubin production, immature liver conjugation, and decreased bilirubin excretion. Pathologic jaundice is defined by onset <24 hours of age, rapid rise (>5 mg/dL/day), elevated direct bilirubin, or persistence beyond 2 weeks
  \item Causes include hemolytic disease, sepsis, Crigler-Najjar syndrome, Gilbert syndrome, biliary atresia, and breast milk jaundice. Unconjugated bilirubin is neurotoxic
  \item At high levels it crosses the blood-brain barrier and causes bilirubin-induced nervous system dysfunction and kernicterus. Diagnosis includes total and direct serum bilirubin, blood type and Coombs test, G6PD screening, and evaluation for sepsis.
\end{itemize}
\section{Risk Factors}
\begin{itemize}[noitemsep]
  \item The right treatment depends on gestational age, risk factors, and bilirubin level: phototherapy is first-line
  \item Exchange transfusion for severe cases or signs of acute bilirubin encephalopathy.
  \item Genetic predisposition and family history
  \item Advancing age
  \item Lifestyle factors (diet, exercise, smoking, alcohol)
  \item Environmental exposures
  \item Underlying medical conditions
  \item Medication use
\end{itemize}
\section{Signs and Symptoms}

\subsection{Common symptoms}
\begin{itemize}[noitemsep]
  \item Symptoms vary depending on the severity and extent of the condition Pain or discomfort in the affected area Functional impairment affecting daily activities Changes in appearance or function of affected tissues Systemic symptoms may include fatigue, fever, or weight changes Symptoms may develop gradually or appear suddenly
  \item Pain or discomfort in the affected area
  \end{itemize}

\subsection{Emergency warning signs}
\begin{itemize}[noitemsep]
  \item Severe or worsening symptoms of neonatal jaundice / hyperbilirubinemia / kernicterus require immediate medical evaluation
\end{itemize}
\section{Progression and Stages}
The condition may progress through different stages of severity over time. Early stages often involve mild or subtle symptoms that may be overlooked. Moderate stages involve more noticeable symptoms affecting daily function. Advanced stages may involve significant impairment and complications.
\section{Complications}
\begin{itemize}[noitemsep]
  \item Disease progression leading to worsening symptoms
  \item Secondary infections or injuries
  \item Side effects of treatment
  \item Reduced quality of life
  \item Emotional and psychological impact
\end{itemize}
\section{Diagnosis}
Comprehensive medical history and physical examination Laboratory tests to assess organ function and rule out other conditions Imaging studies to visualize affected structures Specialized testing depending on the affected system Consultation with relevant specialists Monitoring of disease progression over time

\begin{itemize}[noitemsep]
  \item Comprehensive medical history and physical examination
  \item Laboratory tests to assess organ function
  \item Imaging studies to visualize affected structures
  \item Specialized testing depending on the affected system
  \item Consultation with relevant specialists
\end{itemize}
\section{Prevention}
\begin{itemize}[noitemsep]
  \item Maintain a healthy lifestyle with balanced diet and exercise
  \item Avoid known risk factors
  \item Attend regular medical check-ups for early detection
  \item Follow recommended screening guidelines
\end{itemize}
\section{Treatment Options}
Medications to manage symptoms and address underlying mechanisms Lifestyle modifications and self-care measures Physical therapy and rehabilitation Surgical intervention when indicated Regular monitoring and follow-up care Patient education and support services

\begin{itemize}[noitemsep]
  \item Medications to manage symptoms and address underlying mechanisms
  \item Lifestyle modifications and self-care measures
  \item Physical therapy and rehabilitation
  \item Surgical intervention when indicated
  \item Regular monitoring and follow-up care
\end{itemize}
\section{Living with It}
\begin{itemize}[noitemsep]
  \item Follow your treatment plan as prescribed
  \item Maintain regular follow-up with your healthcare provider
  \item Make necessary lifestyle adjustments
  \item Seek support from family, friends, or support groups
  \item Stay informed about your condition
\end{itemize}
\section{Outlook}
The outlook is generally positive with timely treatment. The outlook varies depending on the specific condition, severity at diagnosis, and response to treatment. Early intervention generally leads to better outcomes. Many conditions can be effectively managed with appropriate treatment. Regular follow-up care is important for monitoring and treatment adjustment.
